\chapter[Template introduction]{How to use the HMI thesis template in practice} % The title in square brackets is the one printed in the Table of contents, and should thus be reasoanbly short
    \label{cha:introduction}
    %
    \onehalfspacing
    %
\epigraph{%.
    The real voyage of discovery consists not in seeking new landscapes, but in having new eyes.}
    {Marcel Proust}
    %
    \begin{enumerate}
        \item In definitions.tex, (re)define the mathematical operators and symbols that you use frequently. Also define other very common abbreviations that you'll need several times, such as chemical formulae, collections of references, etc. 
        %
        \item In packages.tex, define the page geometry with the margins and eventual printing offsets (the default uses the very asthetic french canon des ateliers aspect ratios with 1/3 of the page as margins). Then load any additional \LaTeX packages that you might need.
        %
        \item Review the values in settings.tex, mainly tocdepth and spacing (single, onehalf, double).
        %
        \item In acronyms.tex, define your acronyms paying attention to the use of plural forms. Read up on how to use the acronyms package: \url{https://www.namsu.de/Extra/pakete/Acronym.html}
        %
        \item Define all important mathematical/physical symbols in symbols.tex
        %
        \item Adjust the title page in titlepage.tex. Fulfill all requirements of the faculty to which you submit! Also, provide adequate copyright statements if you are using already published works for your thesis chapters. Always review the requirements of the publishers (ACS, APS, RSC, etc.).
        %
        \item Now, you can start renaming and reordering the chapters defined in IV\_mainmatter using the thesis.tex file, and finally start to work on your thesis!
    \end{enumerate}

    \textbf{General tips for thesis writing and the use of \LaTeX}
    %
    \begin{itemize}
        \item If you recompile the document and get new errors and/or warnings, \textbf{ALWAYS FIX THEM IMMEDIATELY}. Do not let errors accumulate, as it might get close to impossible to debug them once they pile up and start to interact. You will run into compile timeouts if you have a large document with tons of Latex errors.
        %
        \item Do not cramp too much text into one line. This makes debugging very hard, as Latex will only show you "Error in line X", but if line X contains 1000 words, good luck spotting the issue.
        %
        \item Instead, write only one sentence per line. You can even separate two sentences by a \% sign for better visuals
        %
        \item Use indentation to improve the visual clarity
        %
        \item Always label your chapters, sections, figures and tables immediately and in a systematic and understandable way. This will make it a lot easier to use the ref command when you want to refer to the object
        %
        \item Use multiple .tex files for long documents like PhD theses and "glue" them together through the thesis.tex file using the include command. If you put 100000 words in a single .tex file, you'll take forever to find any text that you might need to change
        %
        \item When creating figures, try to use vector formats (.svg, .eps, .pdf in some cases) over raster formats whenever possible. Create your figures exactly according to the page dimensions defined in packages.tex (e.g., a width of 140mm if textwidth=140mm). Alternatively, for many graphs you can also use the tickz package
        %
        \item Read up on the basic rules of good typesetting before you begin. For example, between a value and its corresponding unit, there should be a half space, i.e: The temperature was 500\,K rather than 500 K or 500K. Remember that units and chemical symbols should never be italicized. Also know which kind of dash to use: hyphens (-) for connected words like well-known, En dashes (--) for ranges like 20--30\,mg, and Em dashes (---) for breaking sentences: Birds---such as pigeons---are animals that have wings.
    \end{itemize}